% ============================================================
% Response to Editor -- Forests manuscript forests-4515777
% Prepared: 2026-10-07
% In-place edits 2026-10-12: Figure 3 -> Appendix Figure A3 (station BJS
%   moved to appendix); removed reference to Appendix Figure A2 (station
%   DR figure dropped from paper).
% In-place edits 2026-10-12 (CLCD station): Added CLCD forest cover (100m
%   buffer, ATT=-2.41 pp, SE=2.07) to Comment 2 itemised list; updated
%   "Both estimates" -> "All estimates"; clarified Figure A3 grid does not
%   include CLCD panel.
% In-place edits 2026-10-12 (title): Updated Comment 4 title item to
%   reflect current title including coding-agents subtitle; added sentence
%   noting the coding-agent use is part of the paper's interest.
% In-place edits 2026-10-13 (appendix restructure): A1 buffer-robustness
%   figure split into A1 (BJS ctrl 5km, col 2) and A2 (CS pref, col 5);
%   old A2 DR -> A3; old A3 station BJS -> A4 (Tmax moved to upper left).
%   Comment 1 Figure A1 description updated; Comment 2 A3 -> A4.
% In-place edits 2026-10-13 (appendix reorder): appendix figures reordered
%   to match Table 1 column order: A1=col2 (BJS 5km), A2=col3 (DR 20km),
%   A3=col5 (CS pref), A4=col6 (station).  Comment 1 paragraph updated.
% ============================================================

\documentclass[12pt]{article}
\usepackage[utf8]{inputenc}
\usepackage[T1]{fontenc}
\usepackage[margin=1in]{geometry}
\usepackage{setspace}
\usepackage{natbib}
\bibpunct{(}{)}{;}{a}{,}{,}
\usepackage{hyperref}
\usepackage{enumitem}

\title{Response to Editor\\
       \large Manuscript forests-4515777\\
       \large ``The Effect of China's National Forest City Program on Urban
       Heat Is Not Robustly Identified: A Comment on Ma, Geng, and Zhong
       (2026)''}
\author{Matthew Turner}
\date{October 2026}

\begin{document}

\maketitle

\noindent
I thank the editor for the careful reading of the manuscript and for the
opportunity to revise.  The editor raises five concerns.  I address each
in turn below.  Revised passages in the manuscript are indicated by page
and line number.

% ------------------------------------------------------------------
\section*{Editor Comment 1: Exact Replication on Ma et al.'s Sample}

\noindent\textit{Editor's request:} An exact replication on the
\citet{Ma2026} sample and outcome --- extreme heat days, 243 cities,
2000--2023 --- using heterogeneity-robust estimators.

\bigskip\noindent
I acknowledge that an exact replication on \citeauthor{Ma2026}'s sample
would be the most direct response to their finding.  Unfortunately, the
data required for such a replication are not available to me through any
accessible channel.

\citet{Ma2026} draw their weather station data from the China
Meteorological Administration (CMA), a proprietary network that requires
Chinese institutional affiliation and fees to access.  Their city-level
control variables come from the China City Statistical Yearbook, which is
similarly not available for programmatic download.  I contacted the
authors directly to request their processed city-year panel; they
declined to share it.  As a result, obtaining the 243-city sample used
by \citeauthor{Ma2026} is not feasible without access to the CMA network
or the authors' data.

The one step available to make the analysis more comparable to
\citeauthor{Ma2026} is to adopt their spatial unit.  \citeauthor{Ma2026}
aggregate outcomes at the prefecture level; the original submission used
circular buffers of 5, 10, and 20~km radius.  I have therefore added
prefecture-level boundary extractions to the analysis, using 2005
prefecture boundary polygons (pf05.shp) matched to each city in our
sample.  Of our 275 sample cities, 258 are successfully matched to a
prefecture polygon; the remaining 17 lie in regions (mainly Xinjiang,
Hainan, and parts of Inner Mongolia) not covered by the shapefile and
have missing values for prefecture-level outcomes.  Prefecture-level
results are consistent with the circular-buffer results across all
outcomes: temperature estimates are small, positive, and statistically
insignificant (CS ATT $= +0.055^\circ$C, $SE = 0.087$; DR ATT
$= +0.036^\circ$C, $SE = 0.082$).

The remaining gap in sample coverage --- 99 cities vs.\ \citeauthor{Ma2026}'s
243 --- reflects the difference in station network density between
GHCN-Daily (the publicly available source) and CMA (the proprietary
source used by \citeauthor{Ma2026}).  This limitation is now documented
explicitly in the revised manuscript.

Appendix Figures~A1, A2, and~A3 show event studies for all four satellite
outcomes for the three city-level specifications in columns~(2)--(5) of
Table~1 (excluding column~(1), shown in Figure~1 of the paper).
Figure~A1 shows BJS with controls at the 5~km buffer (column~2);
Figure~A2 shows DR at 20~km (column~3);
Figure~A3 shows CS at the prefecture level (column~5).
All three figures add CLCD forest cover as a fourth panel alongside
LST, NDVI, and tree cover.

% ------------------------------------------------------------------
\section*{Editor Comment 2: Valid Aggregate Inference for the Station
Hot-Days Analysis}

\noindent\textit{Editor's request:} A valid aggregate ATT with a
properly documented confidence interval for the station hot-days outcome.
The previous submission showed the event-study figure but did not report
the aggregate ATT.

\bigskip\noindent
I have added the aggregate ATT for the station hot-days outcome to the
Results section (Section~4, third paragraph).  Table~1, column~(6) now
reports BJS with controls at the station level, using the same Python
implementation as columns~(1)--(2), and Appendix Figure~A4 shows the corresponding
event studies.  BJS with controls values:

\begin{itemize}[nosep]
  \item Hot days ($\geq 35^\circ$C): ATT $= +1.60$ days ($SE = 2.06$;
        95\%~CI $[-2.44,\;+5.64]$).
  \item Summer Tmax: ATT $= +0.114^\circ$C ($SE = 0.184$).
  \item CLCD forest cover (100~m buffer): ATT $= -2.41$~pp ($SE = 2.07$).
\end{itemize}

All estimates are statistically indistinguishable from zero, consistent
with the city-level null results and with the paper's overall conclusion
that an NFCP temperature effect is not robustly identified.  The CLCD
station result corroborates the city-level vegetation null.  Pre-designation
estimates for all station outcomes are consistent with parallel trends
under covariate adjustment.

Appendix Figure~A4 shows BJS with controls event studies for four station-level
outcomes in a $2 \times 2$ grid (LST/Tmax top left, NDVI top right, tree
cover bottom left, hot days bottom right).  CLCD forest cover is reported
in Table~1, column~(6) and the Results text; it is not included as a
separate panel in Figure~A4.

% ------------------------------------------------------------------
\section*{Editor Comment 3: Treatment of Missing Meteorological
Observations}

\noindent\textit{Editor's request:} The paper does not describe how
station-year observations with few valid daily readings are handled.

\bigskip\noindent
I have added the following to the Data and Methods section (Section~3,
station data paragraph):

\begin{quote}
  ``We retain stations within 50~km of a city centroid with at least 30
  valid daily maximum temperature readings per year and average outcomes
  where multiple stations fall within this buffer.  City-years for which
  no station meets this threshold are excluded from the estimation sample.''
\end{quote}

I have also corrected the buffer radius from 20~km to 50~km in this
passage; the earlier text was a typographical error --- the code applies
a 50~km radius, consistent with the station sample size of 99 cities
reported in the same paragraph.

% ------------------------------------------------------------------
\section*{Editor Comment 4: Title and Conclusion}

\noindent\textit{Editor's request:} The title and conclusion are too
strong; ``not robustly identified'' is what the evidence supports, not
``no effect.''

\bigskip\noindent
I agree with this characterisation and have made the following changes.

\begin{enumerate}[nosep]

\item \textbf{Title.} Changed from ``No Effect of China's National Forest
City Program on Urban Heat or Forest Cover: A Comment on Ma, Geng, and
Zhong (2026)'' to ``The Effect of China's National Forest City Program on
Urban Heat Is Not Robustly Identified: Comment on \citet{Ma2026} and
Demonstration of the Use of Coding Agents for Replication Studies.''
I believe part of the interest of this project lies in its use of coding
agents to conduct a replication exercise that would otherwise be
impractical, and the revised title now reflects this.

\item \textbf{Abstract.} Replaced ``find no statistically or economically
significant cooling effect'' with ``find that a cooling effect of NFCP
designation is not robustly identified.''  Replaced ``NFCP designation
does not lead to a decrease in temperature'' with ``our evidence does not
support a robust cooling effect of the NFCP.''

\item \textbf{Discussion (Section~5).} Replaced ``NFCP designation has no
temperature effect'' with ``The temperature effect of NFCP designation is
not robustly identified.''  Added: ``we cannot rule out a modest cooling
effect.''

\item \textbf{Conclusion (Section~6).} Replaced ``The main result in
\citet{Ma2026} does not survive scrutiny'' with ``The main result in
\citet{Ma2026} is not robustly identified.''  Added: ``our evidence is
not sufficient to rule out a modest effect of the NFCP.''

\end{enumerate}

% ------------------------------------------------------------------
\section*{Editor Comment 5: Data and Code Availability; Factual
Inaccuracies About Ma et al.}

\noindent\textit{Editor's request:} Full data and code availability
statement; correction of factual inaccuracies about \citet{Ma2026}.

\bigskip\noindent
I have made the following corrections and additions.

\begin{enumerate}[nosep]

\item \textbf{Data and Code Availability.}  A new section immediately
before the bibliography lists the public sources for each data layer
(MODIS products from NASA EOSDIS, CLCD from Zenodo, GHCN-Daily from
NOAA, NFCP designation lists from the National Forestry and Grassland
Administration) and states that estimation code is available from the
author upon request.

\item \textbf{Pencil-test claim.}  The previous submission stated that
``a straight line fits within the confidence intervals of all pre- and
post-designation coefficients'' in \citeauthor{Ma2026}'s event study.
This is incorrect: the post-designation coefficients in their Figure~1
reach approximately $-8$ extreme heat days by period $+15$, with
relatively narrow confidence intervals by that point.  I have replaced
this passage with a correct characterisation: the pre-designation
estimates are consistent with a mild declining pre-existing trend, and if
such a trend continued post-designation the post-period pattern would be
observationally equivalent to a treatment effect.

\item \textbf{Description of \citeauthor{Ma2026}'s estimator.}  The
previous submission stated that \citeauthor{Ma2026} ``apply this
estimator [BJS] for their event study but not their main results table.''
I have corrected this to: \citeauthor{Ma2026} report TWFE estimates in
their main results table (their Table~5) and use the BJS imputation
estimator only for the event-study pre-trend test in their Figure~1.  The
TWFE estimates, not the BJS estimates, are the basis of their
quantitative conclusion of approximately three fewer extreme heat days per
year.

\end{enumerate}

% ------------------------------------------------------------------
\bibliographystyle{plainnat}
\bibliography{refs}

\end{document}
